% =============================================================================
% 5. DISCUSSION
% =============================================================================
\section{Discussion}
\label{sec:discussion}

The findings reported in Section~\ref{sec:results} are interpreted in this
section against the four Kirkpatrick levels and against the existing
literature. The discussion is organised by level, with each subsection
ending in a synthesis statement that connects the operational finding to
both safety and sustainability outcomes.

\subsection{Reaction: foundations for engagement and efficiency}
\label{subsec:disc-reaction}

The reaction scores observed across the ten sampled sessions
(Fig.~\ref{fig:scores-reaction}) confirm that the redesigned process
preserved participant engagement. While reaction data are sometimes dismissed
as a weak signal \citep{Patel2023}, two observations justify retaining
Level~1 in regulated environments. First, sustained engagement is a
necessary (although not sufficient) condition for the Level~2 and Level~3
outcomes that follow. Second, structured analysis of the items below the
target threshold provides early warning of instructor or content issues,
which avoids the cost of discovering those issues later through poor
behavioural or operational outcomes \citep{ORconnor2024}. The two
sub-9 scores identified in the sampled cohort triggered targeted follow-up
that, in both cases, resolved instructional or scheduling problems before
they propagated into Level~2 or Level~3.

\subsection{Learning: verifying knowledge to prevent future waste}
\label{subsec:disc-learning}

The increase in first-attempt pass rate from 78\% to 86\%
(Fig.~\ref{fig:scores-grades}) is consistent with prior findings that
explicit pass-mark assessment, combined with structured re-training of
those who fail, increases the proportion of participants who reach the
required knowledge level without ever entering production with an
unverified competency \citep{Nakamura2024,Garcia2024}. Importantly, this
benefit is observable only when Level~2 is implemented systematically; the
common shortcut of treating attendance as evidence of competency reproduces
the compliance paradox identified by \citet{Nakamura2024}.

The connection to sustainability is direct: every participant who reaches
the production floor with a verified competency carries a lower probability
of generating internal-failure events that consume rework labour, scrapped
materials and energy.

\subsection{Behaviour: translating knowledge into sustainable practices}
\label{subsec:disc-behaviour}

The 23-percentage-point increase in observed adherence to the work-card
consultation protocol (from 71\% to 94\%) is the most consequential
behavioural change documented in this study. It is consistent with the
cognitive bottleneck argument of \citet{Dismukes2023}: when training is
reinforced with structured on-the-job observation, technicians internalise
the work-card consultation as a default rather than as a discretionary
step. From a Reason-style system-safety perspective \citep{Reason2016,
Leveson2024}, the work-card consultation is a defensive layer that
intercepts errors before they propagate.

The sustainability translation is again direct: every error intercepted at
the level of the work-card avoids the downstream cascade of rework, scrap,
warranty, concession and expedited shipping.

\subsection{Results: quantifying the strategic impact on performance and waste}
\label{subsec:disc-results}

The 28\% reduction in training failures, the 35\% reduction in regulatory
non-conformances, the approximately 6\% turnaround-time improvement and the
across-the-board increase in customer satisfaction
(Fig.~\ref{fig:customer-satisfaction}, Table~\ref{tab:before-after}) jointly
indicate that the Kirkpatrick-based redesign was associated with measurable
movements at the strategic level. The financial translation in
Section~\ref{subsec:results-economic} estimates an annual net saving in
the range of USD~0.75--1.5~M across scenarios (USD~1.19~M base case).

These figures are in the same order of magnitude as those reported in
recent comparable case work in \MRO{} settings \citep{Zhang2025,Sousa2025},
which lends external plausibility despite the single-site limitation and
the absence of a concurrent control group.

\subsection{Synthesis: a framework for integrated compliance and sustainability}
\label{subsec:disc-synthesis}

Taken together, the level-by-level findings support the central proposition
of this study: a Kirkpatrick-based training evaluation framework can be
operationalised in an \AS{}-certified \MRO{} facility in a way that
delivers simultaneously on compliance, safety and sustainability. The
conceptual model in Figure~\ref{fig:conceptual} captures this synthesis:
each Kirkpatrick level is mapped onto a specific cluster of \AS{}
clauses, a specific \SMS{} pillar and a specific class of sustainability
indicators (waste, energy, emissions, social capital). The remainder of
the paper draws managerial and policy implications from this synthesis.
